\section{Two actual orbits following the whole graphic}
\label{sec:lower-main}
Fix $b\in(0,2)$ throughout this section. We use the subfamily
\begin{equation}\label{eq:lower-realizing-family}
 F=y(1+bz-y)+\varepsilon^2(z^2-z+c),\qquad
 G=zy-\nu,\qquad c,\varepsilon>0,\quad \nu\ge0 .
\end{equation}
The weak saddle passage has a power-law exponent below one. We
combine it with a strongly contracting complement whose splitting
changes sign. The complement must be defined even when its output
lies below the stable branch: only its positive outputs enter the
saddle passage, and the boundary between them supplies the first
displacement sign. We therefore construct it on a full signed
rectangle, including the unstable outlet, before choosing periodic
roots.

\subsection{The signed complement}
Set
\[
 z=cX,\qquad y=c\varepsilon^2Y,\qquad
 e=\varepsilon/\sqrt c,\qquad
 \lambda_+=\nu/(c^2\varepsilon^2),\qquad
 \tau=ct_{\rm phys}.
\]
Then the exact equations are
\begin{equation}\label{eq:lower-realizing-scaled}
 X'=e^2\mathcal F(X,Y),\quad Y'=XY-\lambda_+,\quad
 \mathcal F=Y(1+bcX-c^2e^2Y)+1-X+cX^2 .
\end{equation}
Near the smaller positive equilibrium, we select the invariant graph
containing its weak stable branch. The finite-order product construction
of Section~\ref{sec:normal-product}, applied in the proof of
Lemma~\ref{lem:realizing-complement}, gives coordinates $(\xi,\omega)$
and a positive time variable $s$ such that
\begin{equation}\label{eq:lower-exact-power-main}
 \frac{d\xi}{ds}=-r_s\xi,\qquad
 \frac{d\omega}{ds}=\omega,\qquad
 r_s=e^2\beta_s,\qquad 0<\beta_0\le\beta_s\le\beta_1 .
\end{equation}
For each fixed field, $r_s$ is its constant characteristic ratio.
The stable branch is $\omega=0$ and the unstable branch
is $\xi=0$. Fix an incoming cut at $\theta=\theta_b$ and an outgoing
cut
\[
 \Sigma_{\rm out}=\{\omega=\omega_0,\ \xi=a,\ 0\le a\le a_2\}.
\]
The value $a=0$ is the actual unstable outlet. The incoming weak
coordinate $\xi_0=\xi(\theta_b)$ is a positive unit. The local
positive-side passage is exactly
\begin{equation}\label{eq:lower-power-main}
 D_p(s)=\xi_0(p)(s/\omega_0)^{r_s(p)},\qquad 0<s<\omega_0 .
\end{equation}
Here the argument $s$ is the incoming $\omega$-coordinate; at $s=0$
the saddle flight is infinite.

\begin{lemma}\label{lem:realizing-complement}
There are fixed cuts, $a_2,\omega_0>0$, $m_N,K_0>0$, and a
sufficiently small $c_0>0$ with the following property.
The negative cut is $z=-q_g$, with $q_g>0$ fixed.
For each fixed $0<c<c_0$ and all sufficiently small
$\varepsilon>0$, put
\begin{equation}\label{eq:lower-forcing-window}
 \nu_+=4q_gK_0\varepsilon^3e^{-m_N/\varepsilon^2}.
\end{equation}
The forward complement of \eqref{eq:lower-realizing-family}, from
$\Sigma_{\rm out}$ to the incoming cut, exists in finite physical
time on the entire closed rectangle
$[0,a_2]\times[0,\nu_+]$. Its signed output $H_{c,\varepsilon,\nu}(a)$
is continuous there, has a strictly positive source derivative, and
satisfies
\begin{equation}\label{eq:lower-complement-contract}
 |H_p(a)|<\omega_0/4,\qquad
 0<H'_p(a)\le M_\varepsilon:=e^{-m_N/(2\varepsilon^2)} .
\end{equation}
Moreover,
\begin{equation}\label{eq:lower-splitting-signs}
 H_{c,\varepsilon,0}(0)>0,\qquad
 H_{c,\varepsilon,\nu_+}(0)<0 .
\end{equation}
The images of the whole rectangle lie in one fixed neighborhood of
$\Gamma_b$. As $c,e,\lambda_+\to0$, the outgoing
passage and regular fast arc converge to the corresponding fast
pieces of $\Gamma_b$, uniformly in $a$ and $\nu$.
\end{lemma}

Appendix~\ref{sec:lower-complement} proves the lemma by following every
outgoing source through the height-phase exit, parabolic passage,
fast arc, and moving descent. A fixed negative slow interval contracts
both height differences and source derivatives. Its contraction
dominates the subsequent positive amplification and the finite powers
of $\varepsilon^{-1}$ incurred upstream. A final positive-time
adjustment gives the signed output on the incoming cut, including
negative outputs, which are outside the domain of $D_p$. Constants
with a subscript $c$ are fixed after $c$ is chosen and absorbed by
decreasing $\varepsilon$ afterward.

\subsection{A homoclinic boundary and two roots of one field}
Fix $c$ in the lemma's range and a constant $r_0\in(0,1)$.
Before selecting the splitting parameter, choose $\varepsilon$
in that range so that $e^2\beta_1<r_0$. Then $0<r_s(p)<r_0$
holds throughout $0\le\nu\le\nu_+$. Keep $c,\varepsilon$ fixed
and set $\alpha(\nu)=H_{c,\varepsilon,\nu}(0)$. By continuity and
\eqref{eq:lower-splitting-signs},
\[
 \nu_*=\inf\{\nu\in[0,\nu_+]:\alpha(\nu)<0\}
\]
belongs to $(0,\nu_+)$, satisfies $\alpha(\nu_*)=0$, and is
approached by parameters in the negative set. At this parameter
the unstable outlet follows the finite complement to the stable
branch of the same saddle, giving a homoclinic orbit. The sign
change and continuity suffice for this choice of $\nu_*$.

At $\nu=\nu_*$, write $h_1=H'(0)>0$. As $a\downarrow0$,
\begin{equation}\label{eq:homoclinic-positive-displacement}
 \frac{D_p(H_p(a))}{a}
 \sim \xi_0\omega_0^{-r_s}h_1^{r_s}a^{r_s-1}
 \longrightarrow\infty .
\end{equation}
Choose $a_1\in(0,a_2)$ where the return displacement is positive.
On the other hand, $H_p(a_2)\le M_\varepsilon a_2$, so
\begin{equation}\label{eq:lower-right-sign}
 D_p(H_p(a_2))
 \le C_D e^{-\beta_0m_N/(2c)}<a_2/4 .
\end{equation}
The unit bound for $\xi_0$ and the fixed cuts determine $C_D$.
The last inequality is imposed in the initial choice of $c$, before
$\varepsilon$ and $a_1$ are chosen.

Choose a single $\nu>0$ in the negative set sufficiently close
to $\nu_*$ that $H_p(a_1)>0$ and the positive middle displacement
persist. Since $H_p$ is strictly increasing and $H_p(0)<0$,
there is a unique $a_{\min}\in(0,a_1)$ with $H_p(a_{\min})=0$.
The actual returning interval is
\begin{equation}\label{eq:actual-lower-return}
 J_p=(a_{\min},a_2],\qquad
 \Pi_p(a)=D_p(H_p(a)).
\end{equation}
Its left displacement limit is $-a_{\min}<0$. Its middle
displacement at $a_1$ is positive, and its displacement at $a_2$
is negative by \eqref{eq:lower-right-sign}. Thus there are distinct
roots
\begin{equation}\label{eq:two-return-roots}
 a^-\in(a_{\min},a_1),\qquad a^+\in(a_1,a_2),\qquad
 \Pi_p(a^-)=a^-,\quad\Pi_p(a^+)=a^+ .
\end{equation}
For $\alpha<0$ the estimate $H_p(a)\le M_\varepsilon a$
still holds, hence every output of \eqref{eq:actual-lower-return}
lies below $a_2/4$ in the same outgoing collar.
Every root has $H_p(a)>0$ and finite canonical saddle time
$T_s=\log(\omega_0/H_p(a))$.
Since $ds=b_N(\theta)d\tau$ and $b_N\ge b_0>0$, its physical
time is at most $T_s/(cb_0)$.

\subsection{Minimal return and isolation}
The complement and $D_p$ are passages of the same field on the
same cuts. To show that their composition is the first return on
$J_p$, we exclude an earlier intersection with the outgoing cut.
Indeed $\omega$ increases after the outgoing hit until the
horizontal positive cut is reached, then $Y$ increases to $e^{-2}$.
Thereafter the physical height is at least $c^2$ before descent,
well above the $O(c\varepsilon^2)$ outgoing collar.
The regular outer tube is separated from that collar, and the
descending and negative pieces have negative $z$. The signed
connector has $z\le z_v<c\theta_b$. Its incoming adjustment
has $\theta\le\theta_b$, while the outgoing cut has
$\theta\ge\theta_b+\sigma$. Finally $D_p$ increases $\omega$
to $\omega_0$ exactly once. These height and spatial clearances
exclude every earlier hit.

To prove isolation, transfer this return to the physical analytic
section
\[
 S_h=\{y=c\varepsilon^2\delta_1\}
\]
inside the positive outgoing tube, with coordinate $q=-z/c$.
Along the positive local passage $D_p$, put $u=Y-h(X)>0$,
where $h$ is the selected stable graph.
The incoming speed $\mathcal F$ is positive. At a first zero of
$\mathcal F$, the equation for $u=Y-h(X)$ gives $Y'>0$, so
$\mathcal F'=\mathcal F_Y Y'>0$, since $\mathcal F_Y>0$ in the tube.
Thus $\mathcal F>0$ throughout this passage.
At a zero of $Y'=XY-\lambda_+$ its derivative is
$e^2\mathcal F Y>0$, so $Y$ has at most one critical point,
a minimum. Both local endpoint heights are below $\delta_1$;
this leg therefore cannot meet $S_h$ before the intended
continuation beyond $\Sigma_{\rm out}$.
The intended hit has positive vertical speed. The finite map
$L_p(a)=-x_1(a,p)$ from the outgoing cut to $S_h$ is a
$C^j$ orientation-preserving diffeomorphism. Its actual first return is
$R_p=L_p\circ\Pi_p\circ L_p^{-1}$.
At every interior returning point, the polynomial flow and its
transverse hitting time are analytic. Thus $R_p$ is analytic on the
connected physical returning interval. The displacement is not
identically zero because it is positive at $L_p(a_1)$.
Its interior zeros are therefore isolated. The two roots in
\eqref{eq:two-return-roots} give distinct isolated orbits: a single
orbit with two outgoing hits would contradict the established first
return. They lie in $y>0$: on the axis $G=-\nu<0$, so an orbit through
the positive outgoing cut could cross downward only and could not
close after doing so. Simplicity will follow from the upper bound
in Section~\ref{sec:combination}.

\subsection{Each orbit converges to the entire graphic}
Fix an open neighborhood $\mathcal U$ of $\Gamma_b$.
All regular tubes, foot boxes, and axis strips are selected inside it
before the realizing parameters are chosen. The whole complement
rectangle has the following image bounds:
\[
 \begin{array}{ll}
 \text{saddle and initial exit:}&|z|\le Cc,\quad |y|\le Cc\varepsilon^2,\\
 \text{height-phase exit:}&|z|\le Cc,\quad 0<y\le c^2,\\
 \text{positive parabolic passage:}&
             \gamma(c^2+z^2)\le y\le M(c^2+z^2),\\
 \text{negative slow passage:}&
             |y|\le K_0\varepsilon^3+C_h\nu,\\
 \text{signed connector:}&
             |y|\le C_c\varepsilon
                        e^{-(m_N-2I_\sigma c)/\varepsilon^2}.
 \end{array}
\]
The parabolic and outer passages converge on every compact fast
subarc avoiding the endpoints. Their moving descending foot
$\Lambda_p(a)$ converges to $f_b$, uniformly on the source
rectangle. The low slow passage has $F>0$ and covers all
abscissae from $\Lambda_p(a)$ to $z_v=O(c)$.
The final descending tail from $y=\rho$ to $K\varepsilon^3$
has abscissa displacement
$O(\rho)+O(\varepsilon^2(1+|\log\varepsilon|))$,
so it stays in the foot's $O(\rho)+o(1)$ ball.

Choose $c_n\downarrow0$ first, with $c_n<c_{n-1}/2$ and
all fixed $c$ inequalities satisfied. For each $c_n$ choose
$\varepsilon_n$ small enough to satisfy its finite-power smallness conditions,
$e_n<1/n$, $\varepsilon_n<1/n$, the common local windows, and
the uniform image error required at that stage.
Also require $\nu_{+,n}/(c_n^2\varepsilon_n^2)<1/n$.
The joint fast convergence fixes the required smallness of $c_n$
before $\varepsilon_n$ is chosen.
Then choose one $\nu_n$ and both roots as above.

For either root separately, the local and initial parabolic pieces
lie in the origin's $O(\rho)+O(c_n)$ ball until $z=\rho$.
The remaining compact fast piece converges to $\gamma_b$, and the
descending tail lies in the foot's $O(\rho)+o(1)$ ball.
The slow graph has uniformly vanishing height and covers
$[\Lambda_n,z_v]$. These facts give
$\sup_{x\in L_n^\pm}\operatorname{dist}(x,\Gamma_b)\to0$.

For the opposite inclusion, every compact fast subarc avoiding the
endpoint balls is covered by each actual fast passage, with uniform
parametrized convergence. Every point of $[f_b+\rho,-\rho]$
is covered in abscissa by each monotone slow passage, with height
tending to zero. The remaining graphic tails lie in the endpoint
balls; each orbit has a cutoff point tending to the foot and a
slow point at $z=-\rho$. Their distance from those tails is
$O(\rho)+o(1)$. First take $\rho$ small and then $n$ large.
This gives \eqref{eq:two-whole-limits} for each orbit individually
and proves the existence and isolation part of assertion (ii).
